\documentclass[11pt]{article}

\usepackage[margin=1in]{geometry}
\usepackage{amsmath}
\usepackage{booktabs}
\usepackage{hyperref}
\usepackage{listings}
\usepackage{xcolor}
\usepackage{graphicx}

\lstset{
  basicstyle=\ttfamily\footnotesize,
  breaklines=true,
  frame=single,
  columns=fullflexible,
  keepspaces=true
}

\title{Recreating Milo Network Motifs\\
\large with GraphEasy, Python, and mfinder}
\author{GraphEasy Motif Project}
\date{}

\begin{document}
\maketitle
\section{Introduction}

A motif is a small directed pattern (an induced subgraph pattern) that appears in the real network more often than in randomized networks with the same local constraints.

This document recreates and verifies motif and triad-significance computations from:
(1) Milo et al.\ (2002), and
(2) Milo et al.\ (2004).
For both papers we use three independent counters:
GraphEasy, a small Python triad census, and the original mfinder~1.20.
We only use the graphs we can find; the links are provided below.

In this document we completed feed-forward loop (FFL), bi-parallel, three-chain, bi-fan, and other Table~1 motifs; the demo below uses FFL on one graph as an example.

\section{Milo et al.\ 2002: Table~1 (real motif counts)}

On a single original graph (\textit{E.\ coli}), we show FFL as the demo example.
We count the same induced patterns in GraphEasy, Python, and mfinder and require exact equality of real counts $N_{\mathrm{real}}$.

\subsection*{Demo (one graph, FFL)}

\begin{itemize}
\item \textbf{GraphEasy (FFL DSL)}
\begin{lstlisting}
cd milo2002_reproduction

graph G {
  directed: true;
  edges: file "data/grapheasy/ecoli_table1.txt";
};

motifs FeedForwardLoop = [G where motif {
  source->intermediate;
  source->target;
  intermediate->target;
}];

print numMotifs(FeedForwardLoop);
\end{lstlisting}

\item \textbf{Python (exact induced FFL count)}
\begin{lstlisting}
from pathlib import Path

def load_graph(path: Path, vertices: int) -> list[set[int]]:
    outgoing = [set() for _ in range(vertices)]
    seen = set()
    for line in path.read_text().splitlines():
        if not line.strip():
            continue
        source, target = map(int, line.split()[:2])
        edge = (source, target)
        if edge in seen:
            continue
        seen.add(edge)
        outgoing[source].add(target)
    return outgoing

def induced_edges(outgoing, vertices):
    return {
        (u, v) for u in vertices for v in vertices
        if u != v and v in outgoing[u]
    }

def feed_forward_loops(outgoing):
    matches = []
    for source, nbrs in enumerate(outgoing):
        for intermediate in nbrs:
            for target in nbrs & outgoing[intermediate]:
                if len({source, intermediate, target}) != 3:
                    continue
                verts = (source, intermediate, target)
                expected = {
                    (source, intermediate),
                    (source, target),
                    (intermediate, target),
                }
                if induced_edges(outgoing, verts) == expected:
                    matches.append(verts)
    return matches

outgoing = load_graph(Path("data/grapheasy/ecoli_table1.txt"), 424)
print("FFL", len(feed_forward_loops(outgoing)))  # 40
\end{lstlisting}

\item \textbf{mfinder command (real counts)}
\begin{lstlisting}
cd milo2002_reproduction

./reference/mfinder1.2/mfinder \
  data/mfinder/ecoli.txt -s 3 -r 0 -nu \
  -f results/mfinder_ecoli_size3.txt -q
\end{lstlisting}
\end{itemize}

\begin{center}
\begin{tabular}{l l r l}
\toprule
\multicolumn{1}{c}{Network} & \multicolumn{1}{c}{Motif} & \multicolumn{1}{c}{Nreal (GE/Py/mf)} & \multicolumn{1}{c}{Null/Z} \\
\midrule
\textit{E.\ coli} & feed-forward loop (FFL) & 40/40/40 & yes \\
\textit{E.\ coli} & bi-fan & 203/203/203 & yes \\
St.\ Martin & three-chain & 469/469/469 & yes \\
St.\ Martin & bi-parallel & 382/382/382 & yes \\
Skipwith & three-chain & 184/184/184 & yes \\
Skipwith & bi-parallel & 397/397/397 & yes \\
B.\ Brook & three-chain & 181/181/181 & yes \\
B.\ Brook & bi-parallel & 267/267/267 & no$^*$ \\
s208 & three-node feedback loop & 10/10/10 & yes \\
s208 & bi-fan & 4/4/4 & yes \\
s208 & four-node feedback loop & 5/5/5 & yes \\
s420 & three-node feedback loop & 20/20/20 & yes \\
s420 & bi-fan & 10/10/10 & yes \\
s420 & four-node feedback loop & 11/11/11 & yes \\
s838 & three-node feedback loop & 40/40/40 & yes \\
s838 & bi-fan & 22/22/22 & yes \\
s838 & four-node feedback loop & 23/23/23 & yes \\
\bottomrule
\end{tabular}
\end{center}

$^*$Bridge Brook bi-parallel: paper reports $N_{\mathrm{rand}}=30\pm 7$ ($Z=32$),
while our exact triad-preserving ensemble gives $90.0\pm 23.4$ ($Z=7.57$).
The real count is correct in all three tools; only the null statistics disagree.

\section{Milo et al.\ 2004: Superfamilies (separate recreation)}

We use the same three counters (GraphEasy, Python, mfinder) on recovered Alon networks.
For each network we report feed-forward loop (FFL) and three-chain real counts and require exact equality across all three tools.

\begin{center}
\begin{tabular}{lrrrrrr}
\toprule
Network & FFL GE & FFL Py & FFL mf & Chain GE & Chain Py & Chain mf \\
\midrule
TRANSC-E.COLI & 40 & 40 & 40 & 162 & 162 & 162 \\
SOCIAL-1 & 13 & 13 & 13 & 103 & 103 & 103 \\
SOCIAL-3 & 11 & 11 & 11 & 79 & 79 & 79 \\
JAPANESE & 4111 & 4111 & 4111 & 332627 & 332627 & 332627 \\
ENGLISH & 82031 & 82031 & 82031 & 4988512 & 4988512 & 4988512 \\
FRENCH & 12435 & 12435 & 12435 & 2193156 & 2193156 & 2193156 \\
SPANISH & 50790 & 50790 & 50790 & 7479660 & 7479660 & 7479660 \\
\bottomrule
\end{tabular}
\end{center}

\section{Analytic null model (Itzkovitz et al.\ 2003)}

Every count above is a \emph{real} count $N_{\mathrm{real}}$. Calling a subgraph a
motif also needs $N_{\mathrm{rand}}$, and in the sections above that came from
mfinder generating randomized ensembles. That dependency is the weakest link in
the reproduction: it is slow, it is not reproducible (mfinder seeds from the
wall clock), and it is exactly where our Bridge Brook bi-parallel result departs
from the paper.

This section removes it. Itzkovitz, Milo, Kashtan, Ziv and Alon,
\textit{Subgraphs in random networks}, Phys.\ Rev.\ E \textbf{68}, 026127 (2003)
(\url{https://arxiv.org/abs/cond-mat/0302375}) give closed-form expressions for
the expected number of any subgraph in a degree-preserving random ensemble. No
random network is generated at any point; the expectation is a function of the
moments of the real network's degree sequence, so it is deterministic and
$O(N)$.

\subsection{The formulas}

Split the degree sequence into single out-degree $K$, single in-degree $R$, and
mutual degree $M$ (a pair of edges in both directions). For a subgraph with $n$
nodes, $g_a$ single edges and $g_m$ mutual edges, their Eq.~(5) is
\[
\langle G\rangle \;=\;
\frac{a\,N^{\,n-g_a-g_m}}{\langle K\rangle^{g_a}\langle M\rangle^{g_m}}
\prod_{j=1}^{n}
\Big\langle \binom{K_i}{k_j}\binom{R_i}{r_j}\binom{M_i}{m_j}\Big\rangle_i ,
\]
where $\{k_j\},\{r_j\},\{m_j\}$ are the subgraph's own degree sequences and $a$
is a symmetry factor built from $|\mathrm{Aut}|$ --- the same automorphism group
our motif matcher already computes to pick one representative per orbit. For the
feed-forward loop this reduces to
\[
\langle \mathrm{FFL}\rangle=
\frac{\langle K(K-1)\rangle\,\langle RK\rangle\,\langle R(R-1)\rangle}
     {\langle K\rangle^{3}} .
\]

Eq.~(5) counts \emph{non-induced} occurrences: pairs with no specified edge are
left free. Our matcher is induced by construction, so the comparable quantities
are the starred values of their Appendix~B, recovered by inclusion--exclusion,
e.g.\ $\langle \mathrm{id}6^{*}\rangle=\langle \mathrm{id}6\rangle-\langle
\mathrm{id}38\rangle-\langle \mathrm{id}108\rangle$.

\subsection{A fourth, independent counter}

Their Appendix~C gives the thirteen triad counts as adjacency-matrix products.
Writing $M=A+S$ with $A$ the single arrows and $S$ the mutual ones, $\tilde{A}$
the logical inverse and $A'$ the transpose, the feed-forward loop is
$\sum A^{2}\!\cdot\!A$ and the mutual triangle is $\sum (S^{2}\!\cdot\!S)/6$.
This shares no code path with either the GraphEasy matcher or the Python
enumeration --- it is pure linear algebra --- so it is a genuinely independent
check. It is implemented in \texttt{appendixC.py}.

\subsection{Results: the triad census}

We count all thirteen connected three-node directed subgraphs (induced), by
three independent methods: the GraphEasy \texttt{motifs} matcher, a Python
enumeration over connected triples, and the Appendix~C matrix formulas.

\begin{center}
\begin{tabular}{r l rr rr rr}
\toprule
& & \multicolumn{2}{c}{\textit{E.\ coli}} & \multicolumn{2}{c}{\textit{C.\ elegans}} & \multicolumn{2}{c}{s208} \\
\cmidrule(lr){3-4}\cmidrule(lr){5-6}\cmidrule(lr){7-8}
id & subgraph & obs. & exp. & obs. & exp. & obs. & exp. \\
\midrule
6   & out-star            & 226  & 258.6  & 419  & 427.2  & 164 & 161.7 \\
12  & three-chain         & 162  & 194.6  & 872  & 876.7  & 240 & 267.7 \\
14  & out-star + mutual   & 0    & 0      & 91   & 87.2   & 0   & 0     \\
36  & in-star             & 4760 & 4792.6 & 6044 & 6039.3 & 89  & 86.7  \\
38  & feed-forward loop   & \textbf{40} & \textbf{7.4} & 125 & 122.3 & 0 & 2.3 \\
46  & bi-parallel         & 0    & 0      & 2    & 9.3    & 0   & 0     \\
74  & in-star + mutual    & 0    & 0      & 112  & 126.1  & 0   & 0     \\
78  & two mutual (path)   & 0    & 0      & 7    & 6.5    & 0   & 0     \\
98  & three-node feedback & 0    & 0.02   & 0    & 3.0    & \textbf{10} & \textbf{1.0} \\
102 & FFL + mutual        & 0    & 0      & 4    & 2.0    & 0   & 0     \\
108 & one mutual + 2-in   & 0    & 0      & 7    & 1.4    & 0   & 0     \\
110 & two mutual + single & 0    & 0      & 0    & 0.5    & 0   & 0     \\
238 & three mutual        & 0    & 0      & 0    & 0.02   & 0   & 0     \\
\bottomrule
\end{tabular}
\end{center}

All three counters agree on every entry: $13\times3=39$ of $39$ exact.

The scientific reading is the one Milo reports. In \textit{E.\ coli} the
feed-forward loop is the \emph{only} enriched triad, $40$ observed against
$7.4$ expected ($5.4\times$); every degree-driven triad sits at $0.8$--$1.0\times$,
which is the null model behaving as it should. The s208 circuit instead enriches
the three-node feedback loop, $10$ against $1.0$ ($10.3\times$), a different
motif family for a different network class.

\subsection{How the neurons column is computed in GraphEasy}

The whole of Eq.~(5) is arithmetic on the degree sequence, so it runs inside
GraphEasy with no external tool. This subsection shows the program that
produces the numbers validated below, using the neurons network
(\textit{C.\ elegans}, $N=252$, $509$ directed edges, $14$ mutual pairs).

\paragraph{Step 1 --- split the degree sequence.}
$K$ and $R$ count only \emph{single} out- and in-edges; an edge whose reverse
is also present contributes to the \emph{mutual} degree $M$ instead. The
separation matters: the other twelve entries of Table~I only reproduce under
this reading, and taking $K,R$ as total degrees gives id38 $=159.9$ against the
paper's $120$.

\begin{lstlisting}
graph Net {
  directed;
  edges: file "data/grapheasy/c_elegans_table1.txt";
};

int n = numVertices(Net);
int K[n];   // single out-degree
int R[n];   // single in-degree
int M[n];   // mutual degree
int i = 0;
while (i < n) { K[i] = 0; R[i] = 0; M[i] = 0; i = i + 1; }

for each edge u, v in Net {
  if (hasEdge(Net, v, u)) {
    M[u] = M[u] + 1;              // reverse edge exists -> mutual
  } else {
    K[u] = K[u] + 1;
    R[v] = R[v] + 1;
  }
}
\end{lstlisting}

\paragraph{Step 2 --- take the nine moments.}
Every factor in Eq.~(5) is an average over all $N$ vertices of a product of
falling factorials of these three degrees.

\begin{lstlisting}
real sK = 0.0;   real sKK1 = 0.0; real sKR = 0.0;
real sR = 0.0;   real sRR1 = 0.0; real sKM = 0.0;
real sM = 0.0;   real sMM1 = 0.0; real sRM = 0.0;
i = 0;
while (i < n) {
  sK = sK + K[i];   sR = sR + R[i];   sM = sM + M[i];
  sKK1 = sKK1 + K[i] * (K[i] - 1);
  sRR1 = sRR1 + R[i] * (R[i] - 1);
  sMM1 = sMM1 + M[i] * (M[i] - 1);
  sKR = sKR + K[i] * R[i];
  sKM = sKM + K[i] * M[i];
  sRM = sRM + R[i] * M[i];
  i = i + 1;
}
real mK = sK / n;      real mKK1 = sKK1 / n;  real mKR = sKR / n;
real mR = sR / n;      real mRR1 = sRR1 / n;  real mKM = sKM / n;
real mM = sM / n;      real mMM1 = sMM1 / n;  real mRM = sRM / n;
\end{lstlisting}

For the neurons network these evaluate to
\[
\begin{array}{lll}
\langle K\rangle = 1.9087 & \langle R\rangle = 1.9087 & \langle M\rangle = 0.1111\\
\langle K(K-1)\rangle = 4.3730 & \langle R(R-1)\rangle = 48.9762 & \langle M(M-1)\rangle = 0.0556\\
\langle KR\rangle = 3.9722 & \langle KM\rangle = 0.3929 & \langle RM\rangle = 0.5159
\end{array}
\]

\paragraph{Step 3 --- evaluate the thirteen equations.}
Each row of Table~I is one expression in those moments. The feed-forward loop,
for instance, is $\langle K(K-1)\rangle\langle RK\rangle\langle
R(R-1)\rangle/\langle K\rangle^{3}$.

\begin{lstlisting}
real e6   = n * mKK1 / 2.0;                        // out-star
real e12  = n * mKR;                               // three-chain
real e14  = n * mKM;
real e36  = n * mRR1 / 2.0;                        // in-star
real e38  = mKK1 * mKR * mRR1 / dK3;               // feed-forward loop
real e46  = mKM * mKM * mRR1 / (2.0 * dK2 * dM1);
real e74  = n * mRM;
real e78  = n * mMM1 / 2.0;
real e98  = mKR * mKR * mKR / (3.0 * dK3);         // 3-node feedback
real e102 = mKM * mRM * mKR / (dK2 * dM1);
real e108 = mRM * mRM * mKK1 / (2.0 * dK2 * dM1);
real e110 = mKM * mRM * mMM1 / (dK1 * dM2);
real e238 = mMM1 * mMM1 * mMM1 / (6.0 * dM2 * dM1);
\end{lstlisting}

\noindent
\texttt{dK1}\ldots\texttt{dM2} are the powers of $\langle K\rangle$ and
$\langle M\rangle$ appearing in the denominators, guarded against zero: a
network with no mutual edges has $\langle M\rangle=0$, and the six
$M$-dependent rows are then identically zero rather than undefined.

\paragraph{Step 4 --- correct to induced counts.}
Eq.~(5) leaves unmentioned pairs free, so it counts a superset. Appendix~B
subtracts the denser triads that fill those null edges.

\begin{lstlisting}
real i6  = e6  - e38 - e108;
real i12 = e12 - e38 - e102;
real i14 = e14 - e46 - e102 - e110;
real i36 = e36 - e38 - e46;
real i74 = e74 - e102 - e108 - e110;
real i78 = e78 - e110 - e238;
\end{lstlisting}

\noindent
The remaining seven rows need no correction: they are the denser triads that
the corrections above subtract, and have no null edge left to fill.

\paragraph{Step 5 --- report against the observed census.}
The same program declares the thirteen motifs and prints observed and expected
side by side.

\begin{lstlisting}
motifs T38 = [Net where motif { a -> b; b -> c; a -> c; }];
print numMotifs(T38);        // observed  125
print e38;                   // expected  122.339
\end{lstlisting}

\noindent
The complete program is \texttt{milo.graph}: degree decomposition, nine moments,
thirteen equations, six corrections and thirteen motif declarations, with no
external tool and no randomized network.

\subsection{Validation against the paper's own table}

Itzkovitz et al.\ publish their computed values for \textit{C.\ elegans}
(the ``neurons'' column of their Table~I), so our expectations can be checked
directly against the authors rather than only against ourselves. Quoting, as
they do, to two significant figures, and noting their convention that values
below $0.5$ are rounded to zero:

\begin{center}
\begin{tabular}{r r r c}
\toprule
id & paper & ours & \\
\midrule
6 & $4.3\times10^{2}$ & 430 & match \\
12 & $8.7\times10^{2}$ & 880 & match \\
14 & $8.7\times10^{1}$ & 87 & match \\
36 & $6.0\times10^{3}$ & 6000 & match \\
38 & $1.2\times10^{2}$ & 120 & match \\
46 & 9.3 & 9.3 & match \\
74 & $1.3\times10^{2}$ & 130 & match \\
78 & 6.6 & 6.5 & match \\
98 & 4.5 & 3.0 & see below \\
102 & 2 & 2.0 & match \\
108 & 1.4 & 1.4 & match \\
110 & 0 & 0 & match \\
238 & 0 & 0 & match \\
\bottomrule
\end{tabular}
\end{center}

Twelve of thirteen reproduce the published values. The exception is id98, the
three-node feedback loop, where the paper prints $4.5$ and we obtain $3.0$.

We traced this to an arithmetic slip in that one cell. Back-solving the constant
their number requires gives
\[
\frac{\langle KR\rangle^{3}}{c\,\langle K\rangle^{3}}=4.5
\quad\Longrightarrow\quad c=2.003 ,
\]
against the $c=3$ printed in their own formula $\langle KR\rangle^{3}/3\langle
K\rangle^{3}$; the ratio $4.5/3.004=1.498\simeq 3/2$ is exactly this. With
$c=2$ the expression evaluates to $4.5065$, which prints as $4.5$. The degree
moments are not in dispute --- $N=252$, $\langle K\rangle=1.908730$, $\langle
KR\rangle=3.972222$ --- and the other twelve entries confirm both the network
and our reading of $K$ and $R$ as single-edge degrees. (Taking them as total
degrees would give id98 $=5.2$ but also id38 $=159.9$ against their $120$,
so that reading is excluded.)

Two checks internal to the paper establish that $3$ is the correct divisor.
Their own definition of the symmetry factor is $a^{-1}=a_{0}\prod_{j}
k_j!\,r_j!\,m_j!$, and for the directed three-cycle every node has $k=r=1$,
$m=0$, so the product is $1$ and $a_{0}=|\mathrm{Aut}|=3$, the cyclic rotations
$C_3$. Independently, their Appendix~C formula for the same subgraph is
$\sum (A'^{2}\!\cdot\!A)/3$, also dividing by three. The paper therefore
specifies $3$ in two places and prints a value consistent with $2$ in a third.
We conclude our $3.0$ is correct and the published $4.5$ is a slip. It is
visible only in the neurons column: for transcription the value is $0.02$,
which rounds to zero under either divisor, and we do not have their www network.

\subsection{What this changes}

The randomization machinery is no longer required for the mean. Where the
sections above needed mfinder to estimate $N_{\mathrm{rand}}$, the expectation is
now a closed form in the degree moments, computed inside GraphEasy itself
(\texttt{milo.graph}: degree decomposition, nine moments, thirteen Table~I
formulas, and the Appendix~B corrections, in one program with no external tool).

One limitation is worth stating plainly. The paper derives means only; it gives
no standard deviation, so a $Z$-score in the strict sense is still unavailable
from this method alone, and the ratios above are enrichment factors, not
significance values. Establishing $\sigma$ still requires either a small
randomized sample or a later result. What has been removed is the expensive and
irreproducible part --- the mean over a thousand randomized networks --- not the
whole of the statistics. The approximation also assumes sparsity
($\langle K\rangle\ll N$); \textit{E.\ coli} ($\langle K\rangle=1.22$) sits well
inside that regime, \textit{C.\ elegans} ($1.91$) less comfortably, which is
consistent with the deviations the authors themselves report for the neuron
network.

\section{Source links}

\subsection{Milo et al.\ 2002 inputs (Table~1)}

\begin{itemize}
  \item \textit{E.\ coli}: \url{https://web.archive.org/web/20030116052450id_/http://www.weizmann.ac.il/mcb/UriAlon/Network_motifs_in_coli/ColiNet-1.0/coliInterNoAutoRegVec.txt}
  \item Circuits \textit{s208/s420/s838}: \url{https://web.archive.org/web/20140502072605id_/http://www.weizmann.ac.il:80/mcb/UriAlon/sites/mcb.UriAlon/files/uploads/CollectionsOfComplexNetwroks/s208_st.txt}
  \item Food webs (installer containing WDFs): \url{https://web.archive.org/web/20060502091624id_/http://www.foodwebs.org:80/download/FoodWeb3DInstall.exe}
\end{itemize}

\subsection{Milo et al.\ 2004 inputs (Superfamilies)}

\begin{itemize}
  \item TRANSC-E.COLI (\texttt{coliInterNoAutoRegVec.txt}): \url{https://web.archive.org/web/20030116052450id_/http://www.weizmann.ac.il/mcb/UriAlon/Network_motifs_in_coli/ColiNet-1.0/coliInterNoAutoRegVec.txt}
  \item SOCIAL-1 (\texttt{prisoninter\_st.txt}): \url{https://web.archive.org/web/20140502072605id_/http://www.weizmann.ac.il:80/mcb/UriAlon/sites/mcb.UriAlon/files/uploads/CollectionsOfComplexNetwroks/prisoninter_st.txt}
  \item SOCIAL-3 (\texttt{leader2inter\_st.txt}): \url{https://web.archive.org/web/20140502072605id_/http://www.weizmann.ac.il:80/mcb/UriAlon/sites/mcb.UriAlon/files/uploads/CollectionsOfComplexNetwroks/leader2inter_st.txt}
  \item ENGLISH (\texttt{darwinbookinter\_st.txt}): \url{https://web.archive.org/web/20140501000000id_/http://www.weizmann.ac.il/mcb/UriAlon/sites/mcb.UriAlon/files/uploads/CollectionsOfComplexNetwroks/darwinbookinter_st.txt}
  \item FRENCH (\texttt{frenchbookinter\_st.txt}): \url{https://web.archive.org/web/20140502072605id_/http://www.weizmann.ac.il:80/mcb/UriAlon/sites/mcb.UriAlon/files/uploads/CollectionsOfComplexNetwroks/frenchbookinter_st.txt}
  \item SPANISH (\texttt{spanishbookinter\_st.txt}): \url{https://web.archive.org/web/20140502072605id_/http://www.weizmann.ac.il:80/mcb/UriAlon/sites/mcb.UriAlon/files/uploads/CollectionsOfComplexNetwroks/spanishbookinter_st.txt}
  \item JAPANESE (\texttt{japanesebookinter\_st.txt}): \url{https://web.archive.org/web/20140502072605id_/http://www.weizmann.ac.il:80/mcb/UriAlon/sites/mcb.UriAlon/files/uploads/CollectionsOfComplexNetwroks/japanesebookinter_st.txt}
\end{itemize}

\end{document}
